% Generated draft: requires mathematical and source review.
\documentclass[11pt]{article}
\usepackage[margin=25mm]{geometry}
\usepackage{fontspec}
\setmainfont{Latin Modern Roman}
\usepackage{amsmath,amssymb,mathtools,unicode-math}
\setmathfont{Latin Modern Math}
\usepackage{enumitem,needspace,xurl,hyperref}
\setlength{\parindent}{0pt}
\setlength{\parskip}{5pt}
\setlength{\emergencystretch}{3em}
\newcommand{\N}{\mathbb{N}}
\newcommand{\Z}{\mathbb{Z}}
\newcommand{\Q}{\mathbb{Q}}
\newcommand{\R}{\mathbb{R}}
\newcommand{\C}{\mathbb{C}}
\newcommand{\F}{\mathbb{F}}
\newcommand{\A}{\mathbb{A}}
\newcommand{\PP}{\mathbb{P}}
\newcommand{\E}{\mathbb{E}}
\newcommand{\eps}{\varepsilon}
\newcommand{\dd}{\,\mathrm{d}}
\begin{document}
\section*{UnsolvedMath ID: 2306026}
\textbf{Source problem code:} catalogue entry
\subsection{Definitions and mathematical statement}
Let $p>0$ and let $f(z)=\sum_{n\ge0}a_nz^n$ be holomorphic and zero-free in $\mathbb D$, so $a_0\ne0$. In the usual circumferential mean-valency convention, if $n_f(w)$ counts preimages of $w$ in $\mathbb D$ with multiplicity, the condition is
\[
\frac1{2\pi}\int_0^{2\pi}n_f(Re^{i\theta})\,\dd\theta\le p\qquad(R>0),
\]
with the appropriate almost-everywhere interpretation. The source says zero-freeness follows from this condition when $p<1$, and assumes it separately in the general question. Its exact mean-valency convention should be confirmed.

The proposed coefficient bound, at least for $p\ge1$, is printed as $|a_n|\le A_{n,p}$, with comparison function
\[
F(z)=a_0\left(\frac{1+z}{1-z}\right)^{2p}
=\sum_{n\ge0}A_{n,p}z^{2n}.
\]
This literal expansion is inconsistent for $p>0$ and $a_0\ne0$: its right side has only even powers, whereas the displayed function need not. Moreover $a_0$ is not stated to be real positive, so the $A_{n,p}$ as printed need not be real numbers against which an absolute value can be ordered. Neither issue is silently repaired.

A natural normalized interpretation, requiring confirmation, defines real comparison coefficients by
\[
\left(\frac{1+z}{1-z}\right)^{2p}=\sum_{n\ge0}c_n(p)z^n,
\]
using the holomorphic branch equal to $1$ at $z=0$, and asks whether
\[
|a_n|\le |a_0|c_n(p)\qquad(n\ge1)
\]
for all zero-free circumferentially mean-$p$-valent $f$ when $p\ge1$. This is explicitly a candidate reconstruction of the printed inequality, not an asserted corrected transcription.

The source reports that its intended coefficient conjecture holds for $n=1$ for all $p$, and for $n=2,3$ when $p=1$; fails for small positive $p$ and large $n$; and, for each fixed $f$, holds eventually in $n$ if $p>1/4$. Its small-$p$ discussion invokes sufficiently large translates of bounded univalent functions and the unjustified coefficient decay that a universal small-$p$ bound would force. In the univalent $p=1$ subcase it gives the Littlewood bound
\[
|a_n|\le4|a_0|n,
\]
and reports it proved through the Bieberbach/de Branges result and $K_0=1$ from Problem 6.2. This univalent subcase is distinct from the full mean-$1$-valent class.

A second proposed comparison is $I_\lambda(r,f)\le I_\lambda(r,F)$ for $\lambda p>1$ and $0\le r<1$. The extract does not define $I_\lambda$; under the usual positive integral-mean interpretation it compares angular means of $|f|^\lambda$ and $|F|^\lambda$. This convention, the corrected comparison normalization, and the range of $p$ intended for the second question remain to be fixed. All reported special cases and status claims are source-attributed, not independently verified.
\subsection{Short English statement}
Bound coefficients and integral means of zero-free circumferentially mean-valent disk functions by a stated extremal comparison function. The printed comparison expansion has an even-power error and lacks a positive coefficient normalization. The source reports the univalent mean-one case solved, while the broader coefficient and integral-mean questions remain qualified and need their exact conventions restored.
\subsection{Sources}
\begin{itemize}
\item Local source registry: \url{data/unsolvedmath/problems.json}, permanent UnsolvedMath ID $2306026$, source code catalogue entry, including the supplied PARTIAL-PROGRESS summary. Dataset attribution: Ulam AI and UnsolvedMath Contributors; original-source attribution and rights are retained.
\item Hayman, \emph{Research Problems in Function Theory}, source-list date 2018, Problem 6.26: \url{https://arxiv.org/abs/1809.07200}. Source labels 6.10 and 6.11 identify the printed inequality and comparison function; Problems 6.25, 6.2, and Update 6.2 are internal references. The metadata reports source-TeX extraction and July 29, 2026 access, with status and rights review flags.
\item The statement cites Hayman for the eventual-in-$n$ result and Nehari and Bombieri for the $4K_0|a_0|n$ univalent estimate, but the separate references are missing. The supplied summary attributes the resolved univalent case to de Branges via $K_0=1$. No complete additional bibliography or linked-source inspection is supplied here.
\end{itemize}
\end{document}
